\section{The conjecture on Grothendieck pairs}\label{sec:grothendieck}

We now apply the characterization to the Grothendieck-pair conjecture.
The deduction through complemented Euclidean subspaces is the one in
\cite[Corollary 4.8]{GMR25}; the main theorem allows ordinary AP here.
For an algebraic tensor \(t\in E\otimes F\), recall the injective
and projective norms
\[
\|t\|_\eps
=\sup_{\substack{x^*\in B_{E^*}\\y^*\in B_{F^*}}}
  |(x^*\otimes y^*)(t)|,\qquad
\|t\|_\pi
=\inf_{t=\sum_i x_i\otimes y_i}\sum_i\|x_i\|\|y_i\|.
\]
Their completions are denoted by
\(E\widehat\otimes_\eps F\) and \(E\widehat\otimes_\pi F\).
Both norms satisfy the metric mapping property:
tensoring bounded maps multiplies the norm by at most the product
of their operator norms. This follows directly from the two
definitions.

Let \(\mathcal E=(E_n)_{n\ge1}\) be a sequence of real Banach spaces
with \(\dim E_n=n\).
Following \cite[Definition 1.2]{GMR25}, call \((\mathcal E,F)\)
a \emph{Grothendieck pair} if there is \(C<\infty\) such that
\begin{equation}\label{eq:grothendieck-pair}
\|A\otimes\Id_F:
E_m\widehat\otimes_\eps F
\longrightarrow E_n^*\widehat\otimes_\pi F\|
\le C\|A\|
\end{equation}
for all \(m,n\ge1\) and every \(A:E_m\to E_n^*\).
The existence of the bounded extension in
\eqref{eq:grothendieck-pair} is part of the condition.

\begin{corollary}\label{cor:grothendieck}
If \((\mathcal E,F)\) is a Grothendieck pair and \(F\) is
infinite-dimensional with the approximation property, then
\((\mathcal E,\ell_2)\) is a Grothendieck pair.
In particular, Pisier's conjecture
\cite[Conjecture 1.4]{GMR25} holds, with the approximation property
in place of the bounded approximation property.
\end{corollary}
\begin{proof}
By Gupta--Misra--Ray \cite[Theorem 1.3]{GMR25}, the pair condition
implies finite cotype of \(F\) and \(F^*\).
Thus Theorem~\ref{thm:main} makes \(F\) \(K\)-convex.

We use the complemented Euclidean subspace theorem of
Figiel--Tomczak-Jaegermann \cite{FTJ79}, in the formulation
\cite[Theorem~4]{Pis84ICM}: for an infinite-dimensional
\(K\)-convex space \(F\), there is \(D<\infty\) such that for every
\(d\ge1\) there are bounded maps
\[
U_d:\ell_2^d\to F,\qquad V_d:F\to\ell_2^d,\qquad
V_dU_d=\Id_{\ell_2^d},\qquad \|U_d\|\|V_d\|\le D.
\]
By the metric mapping property and the factorization
\[
A\otimes\Id_{\ell_2^d}
=(\Id_{E_n^*}\otimes V_d)
 (A\otimes\Id_F)
 (\Id_{E_m}\otimes U_d),
\]
the corresponding tensor operator with \(\ell_2^d\) in place of
\(F\) has norm at most \(CD\|A\|\), uniformly in \(d,m,n,A\).

Every algebraic tensor in \(E_m\otimes\ell_2\) has its second
coordinates in a finite-dimensional Hilbert subspace \(H\subset\ell_2\).
The inclusion of \(H\) and the orthogonal projection onto \(H\)
both have norm one. The metric mapping property therefore shows
that both tensor norms on tensors supported in \(H\) agree with
the corresponding norms computed in \(\ell_2\).
The uniform estimate for finite-dimensional Hilbert spaces
consequently holds on \(E_m\otimes\ell_2\), and extends to its
injective completion. This is \eqref{eq:grothendieck-pair} for
\(\ell_2\), with constant \(CD\).
\end{proof}
